\appendix
\renewcommand{\appendixname}{APPENDIX}
\captionsetup[table]{labelfont=bf,name=Table\space,labelsep=colon,justification=centering}
\captionsetup[figure]{labelfont=bf,name=Figure\space,labelsep=colon,justification=centering}
\renewcommand\thefigure{\thechapter\arabic{figure}} 
\newcommand{\stoptocwriting}{%
  \addtocontents{toc}{\protect\setcounter{tocdepth}{-5}}}
\newcommand{\resumetocwriting}{%
  \addtocontents{toc}{\protect\setcounter{tocdepth}{\arabic{tocdepth}}}}

\captionsetup{
  textfont=bf,
  justification=centering
}

\newpage
\phantomsection
\addcontentsline{toc}{chapter}{\bf APPENDICES}
\begin{center}
    \vspace*{\fill}
    \large{\bf APPENDICES}
    \vspace*{\fill}
\end{center}
\newpage

\phantomsection
\addcontentsline{toc}{section}{\bf {APPENDIX A: LSTM MODEL IMPLEMENTATION}}
\stoptocwriting
\chapter{LSTM MODEL IMPLEMENTATION}
\resumetocwriting

The following PyTorch code demonstrates the implementation of the LSTM model used in this study.

\begin{lstlisting}[language=Python]
import torch
import torch.nn as nn

class LSTMModel(nn.Module):
    def __init__(self, input_size, hidden_size=128, num_layers=1, dropout=0.3):
        super(LSTMModel, self).__init__()
        self.hidden_size = hidden_size
        self.num_layers = num_layers
        
        self.lstm = nn.LSTM(input_size, hidden_size, num_layers, 
                           batch_first=True, dropout=dropout if num_layers > 1 else 0)
        self.dropout = nn.Dropout(dropout)
        self.fc = nn.Linear(hidden_size, 1)
        
    def forward(self, x):
        # x shape: (batch_size, sequence_length, input_size)
        lstm_out, _ = self.lstm(x)
        
        # Take output from the last time step
        last_output = lstm_out[:, -1, :]
        
        # Apply dropout and fully connected layer
        last_output = self.dropout(last_output)
        output = self.fc(last_output)
        return output
\end{lstlisting}

\newpage
\phantomsection
\addcontentsline{toc}{section}{\bf {APPENDIX B: FEATURE ENGINEERING}}
\stoptocwriting
\chapter{FEATURE ENGINEERING}
\resumetocwriting

The following features were engineered from the raw data to capture temporal patterns and hydrological dynamics.

\section*{1. Time Features}
\begin{itemize}
    \item \texttt{hour}: Hour of the day (0-23)
    \item \texttt{day\_of\_week}: Day of the week (0-6)
    \item \texttt{month}: Month of the year (1-12)
    \item \texttt{cyclical\_encoding}: Sine and cosine transformations of hour, day, and month to preserve cyclical nature.
\end{itemize}

\section*{2. Lag Features}
Past values of water level were used as predictors:
\begin{itemize}
    \item \texttt{water\_level\_lag\_1}: Value 1 hour ago
    \item \texttt{water\_level\_lag\_6}: Value 6 hours ago
    \item \texttt{water\_level\_lag\_12}: Value 12 hours ago
    \item \texttt{water\_level\_lag\_24}: Value 24 hours ago
\end{itemize}

\section*{3. Rolling Statistics}
Rolling window statistics were calculated to capture trends:
\begin{itemize}
    \item \texttt{rolling\_mean\_[6,12,24]}: Moving average
    \item \texttt{rolling\_std\_[6,12,24]}: Moving standard deviation (volatility)
    \item \texttt{rolling\_min/max\_[6,12,24]}: Range of values
\end{itemize}

\newpage
\section*{4. Risk Classification Logic}
The function below defines the risk levels based on water level percentages:
\begin{lstlisting}[language=Python, caption=Risk Classification Logic]
# Adapted from ThaiWater (HII) standards
def determine_risk_level(water_level_pct):
    if water_level_pct >= 90:
        return 'Critical'
    elif water_level_pct >= 70:
        return 'High Risk'
    elif water_level_pct >= 30:
        return 'Medium Risk'
    else:
        return 'Low Risk'
\end{lstlisting}

\newpage
\phantomsection
\addcontentsline{toc}{section}{\bf {APPENDIX C: AUTHOR CONTACT INFORMATION}}
\stoptocwriting
\chapter{AUTHOR CONTACT INFORMATION}
\resumetocwriting

\begin{itemize}
    \item \textbf{Dechathon Niamsaard}
    \begin{itemize}
        \item Student ID: 126235
        \item Email: st126235@ait.asia
    \end{itemize}
    
    \item \textbf{Zwe Yu Ya Kyaw Zin Oo}
    \begin{itemize}
        \item Student ID: 125990
        \item Email: st125990@ait.asia
    \end{itemize}
\end{itemize}
